\section{Cronograma}

\pendente{Ajustar as etapas e os períodos ao calendário acadêmico vigente e
validar com a orientadora. Esta tabela é a fonte canônica do cronograma: as
datas reais de execução ficam nos registros de sessão do repositório, e nunca
o contrário.}

\begin{table}[H]
  \centering
  \caption{Cronograma previsto de execução}
  \label{tab:cronograma}
  \begin{singlespace}
  \small
  \begin{tabularx}{\textwidth}{@{}lXl@{}}
    \toprule
    \textbf{Etapa} & \textbf{Atividade} & \textbf{Período} \\
    \midrule
    1 & Definição do tema e do problema de pesquisa & 2026.2 \\
    2 & Revisão bibliográfica e trabalhos relacionados & 2026.2 \\
    3 & Levantamento de requisitos com o setor de licitações & 2026.2 \\
    4 & Modelagem de dados e definição da arquitetura & 2026.2 \\
    5 & Qualificação do projeto (TCC I) & fim de 2026.2 \\
    6 & Implementação do sistema & 2027.1 \\
    7 & Avaliação com usuários & 2027.1 \\
    8 & Redação final e defesa (TCC II) & fim de 2027.1 \\
    \bottomrule
  \end{tabularx}
  \end{singlespace}
  \fonte{Elaborado pelo autor.}
\end{table}
